% ============================================================
% 08discusion.tex
% Discusión de resultados, hallazgos de XAI y cumplimiento de objetivos
% ============================================================

\section{Discusión}
\label{sec:discusion}

\subsection{Respuesta a la pregunta de investigación}
\label{subsec:respuestapi}

La respuesta a la pregunta de la Sección~\ref{subsec:pregunta} es
afirmativa: la probabilidad de victoria sí puede estimarse a partir del
estado observable. B4 supera a la heurística de dominio (B1) en los tres
escenarios, su AUC crece con la partida y nunca cae al nivel del azar, ni
siquiera con días posteriores o arquetipos nuevos, donde se mantiene cerca
de 0.79. La calidad de la estimación depende de la etapa: es modesta al
inicio (AUC de 0.62 a 0.68) y alta al final (más de 0.91).

Las referencias responden las dos preguntas del \textit{track}. B4 supera a
B0 por casi 0.33 de AUC, así que el modelo aprende. B4 supera a B2 por
0.036 a 0.071, así que la complejidad se justifica: el estado tiene
interacciones que una regresión logística no capta.

\subsection{Lectura de las hipótesis}
\label{subsec:discusionhipotesis}

\noindent\textbf{H1, sostenida.} El estado predice el resultado desde el
inicio, y la predicción mejora con la partida. La mejora también aparece en
las partidas que llegan a las tres etapas, así que no se explica porque las
partidas largas sean más fáciles. SHAP muestra que el progreso pesa solo el
6.3\,\%: el modelo no mejora por contar turnos, sino porque el tablero
acumula diferencias de premios, descartes y mazo.

\noindent\textbf{H2, no sostenida en general.} En A, al inicio, el modelo
que solo conoce el turno, quién empezó y qué Pok\'emon están activos
predice igual que B4. Con la regla del protocolo, H2 queda refutada en el
escenario principal. En B y C el resultado cambia: el modelo restringido
cae a 0.589 y 0.535, mientras B4 se mantiene en 0.630 y 0.623. Conocer los
Pok\'emon activos solo sirve cuando la población se repite; leer el estado
sirve también con días y arquetipos nuevos. Bertram \textit{et al.}
describen la misma diferencia entre identificar cartas y describirlas
\cite{bertram_cardrepresentations_2024}. Este resultado confirma que fue
correcto dejar los identificadores de carta fuera del modelo principal.

\noindent\textbf{H3, refutada.} Guo \textit{et al.} y Xenopoulos \textit{et
al.} observaron que, cuando cambia la población, la calibración empeora más
que el ordenamiento \cite{guo_temporalshift_2021,xenopoulos_winprob_2022}.
Aquí ocurre lo contrario. Una explicación posible es que los mazos y los
agentes cambian qué posiciones son buenas, no con qué frecuencia se gana.
El resultado depende de cómo se mida: comparar porcentajes de caída habría
apoyado H3. La descomposición en calibración y ordenamiento permite
comparar ambas pérdidas en las mismas unidades.

\subsection{Hallazgos de explicabilidad}
\label{subsec:discusionxai}

\noindent\textbf{¿Qué variables dominan?} Las de recursos: premios,
descartes y mazo. Recursos y tablero explican el 72.8\,\% de la predicción,
y las ablaciones lo confirman, porque es la única familia cuya eliminación
baja el AUC de forma clara.

\noindent\textbf{¿Tiene sentido que dominen?} Sí. Los premios son la
condición de victoria. Los descartes cuentan el material perdido: muchos
descartes propios bajan la probabilidad y muchos del rival la suben. La
señal más fuerte aparece cuando al rival le quedan pocas cartas en el mazo
(hasta $+6$ en logit). Tiene sentido porque en este corpus casi ninguna
partida termina por premios, así que dejar al rival sin mazo es una forma
común de ganar.

\noindent\textbf{¿Revela algún problema en los datos?} No hay señales de
fuga: todas las variables dominantes son información pública del tablero.
Sí aparece una carencia de la representación. Un activo rival con muchos
puntos de vida máximos sube la probabilidad propia. Una explicación posible
es que los Pok\'emon ``ex'' entregan dos premios al ser derrotados y solo se
reconocen por sus puntos de vida altos. Es una hipótesis, no un resultado
verificado.

De estos hallazgos salen tres acciones concretas:
\begin{enumerate}
  \item Agregar una característica con los premios que entrega cada
        Pok\'emon activo, para que el modelo no dependa de los puntos de
        vida como sustituto.
  \item Vigilar en el agente que el evaluador no sobrevalore dejar al rival
        sin mazo. Es una estrategia válida, pero debe comprobarse.
  \item Explicar las predicciones del agente en términos de premios,
        descartes y mazo, que son las variables con que decide el modelo.
\end{enumerate}

\subsection{Cumplimiento de los objetivos}
\label{subsec:objetivoscumplidos}

La Tabla~\ref{tab:objetivos} relaciona cada objetivo específico de la
Sección~\ref{subsec:objetivos} con su evidencia. Se cumplieron los seis. El
tercero terminó en un resultado contrario a H3, pero se cumplió, porque su
propósito era contrastar la hipótesis, no confirmarla.

\begin{table}[!t]
\centering
\caption{Objetivos específicos y evidencia de su cumplimiento.}
\label{tab:objetivos}
\footnotesize
\begin{tabularx}{\columnwidth}{@{}P{2.2cm} Y@{}}
\toprule
\textbf{Objetivo} & \textbf{Evidencia}\\
\midrule
1. Evolución de la estimación & AUC por etapa y por turno, con el número
de partidas de cada medición (Tabla~\ref{tab:h1}, Fig.~\ref{fig:aucturno})\\
2. Generalización estructural & AUC de 0.833 en partidas nuevas (A) y de
0.787 en arquetipos nuevos (C)\\
3. Estabilidad temporal & AUC de 0.789 en días posteriores (B); H3
contrastada (Tabla~\ref{tab:h3})\\
4. Formalización & Clasificación binaria probabilística con la partida
como unidad independiente (Sección~\ref{sec:metodologia})\\
5. Representación auditada & Tabla de 55\,196 partidas y 62
características, con cobertura de cartas medida
(Sección~\ref{sec:analisistecnico})\\
6. Escalera de referencias & B0 $<$ B1 $<$ B2 $<$ B3 $<$ B4 en los tres
escenarios, 24 pruebas de fuga, SHAP y ablaciones\\
\bottomrule
\end{tabularx}
\end{table}

\subsection{Limitaciones}
\label{subsec:limitaciones}

\begin{itemize}
  \item Se usó una sola semilla, y los intervalos no incluyen la
        variación de la búsqueda de hiperparámetros. Pueden ser algo
        optimistas \cite{brill_winprob_2026}.
  \item En B, el cambio de fecha y el cambio de versión del juego ocurren
        juntos y no pueden separarse.
  \item Los datos son partidas entre agentes automáticos durante diez
        días. No representan el juego entre personas.
  \item El modelo no cubre los estados de preparación ni los ya decididos.
  \item Los Pok\'emon activos solo aproximan la identidad del mazo, lo que
        limita la lectura de H2.
\end{itemize}